% sec-10: what is not claimed, and what would falsify the case.  Figure 20.
\section{Risks and Limits}

\draftinstr{full-apply writes 4,000 to 5,000 characters. This section must be
written against the paper, not for it. Source the limitations from
\appfile{funding/supplementary/source-files/Daraxonrasib-Efficient-LLM-Trial-Simulations.zip}
(the authors' own stated limitations), from
\appfile{funding/potential-partners/UC-San-Diego/priority-steps.md} \S1 (the
three positioning corrections), and from the back matter of the ten
applications.}

\subsection{What Is Not Claimed}

\draftinstr{full-apply builds Table 13: columns Claim a reader might infer, and
What is actually true. At least six rows, including: that the simulations
validate RASolute 302; that daraxonrasib is first-in-human; that a robotic
configuration has been specified or cleared; that drug supply or a
cross-reference exists; that UC San Diego has agreed to anything; that any
patient has been treated.}

\subsection{What Would Falsify the Case}

\figslot{graphviz-type / record nodes}{5.6}{Five verification artifacts and the
reviewer question each field answers. Every field is released with the cohort,
and the two marked fields are the ones that can falsify the programme's central
claim.}

\draftinstr{Two paragraphs naming the two falsifying fields explicitly:
\texttt{credibility\_score} failing to reproduce at 81.9, and
\texttt{negative\_results} being absent from a released cohort. State that both
are released whatever they show.}

\subsection{Risks the Programme Cannot Retire Alone}

\draftinstr{full-apply builds Table 14: columns Risk, Who can retire it, and
What the applicant can do meanwhile. Rows: no qualified site; no drug supply; no
sponsor-investigator determination; robotic configuration not permissible;
accrual below plan.}
